\section{ChainRAG}

\begin{frame}{ChainRAG: Metadata}
\small
\textbf{ChainRAG: Mitigating Lost-in-Retrieval Problems in Retrieval Augmented Multi-hop QA}

\paperinfo
{ACL 2025 long paper: Proceedings of ACL 2025}
{OpenAI \texttt{text-embedding-3-small}; BGE-Reranker cross-encoder}
{GPT-4o-mini, Qwen2.5-72B, and GLM-4-Plus}
{Progressive retrieval and rewriting for multi-hop QA}

\begin{block}{Key insights / contributions}
\scriptsize
\begin{itemize}
  \item Formalizes ``lost-in-retrieval'', where missing entities in sub-questions break multi-hop retrieval.
  \item Uses progressive sentence-graph retrieval plus sub-question rewriting to recover missing bridge entities.
  \item Improves effectiveness and efficiency across multiple MHQA datasets and reader LLMs.
\end{itemize}
\end{block}
\end{frame}

\begin{frame}{ChainRAG: Problem and Failure Mode}
\scriptsize
\begin{columns}[T,onlytextwidth]
  \begin{column}{0.48\textwidth}
    \begin{block}{Problem}
      Question decomposition helps MHQA, but later sub-questions often contain pronouns or vague references.
    \end{block}
    \begin{block}{Empirical failure}
      The paper reports that the second sub-question has much lower Recall@2 than the first one, with an average decrease of 18.29\% across three datasets.
    \end{block}
  \end{column}
  \begin{column}{0.48\textwidth}
    \begin{block}{Why existing RAG fails}
      Naive decomposition retrieves with "this author" or "this place" instead of the resolved entity, so retrieval follows a wrong branch.
    \end{block}
    \begin{block}{Method response}
      ChainRAG answers each hop, rewrites the next sub-question with the recovered entity, and retrieves over a sentence graph.
    \end{block}
  \end{column}
\end{columns}
\end{frame}

\pipelineframe{ChainRAG}{images/chain_rag_pipeline.png}{0.98}

\begin{frame}{ChainRAG: Algorithm Deep Dive}
\scriptsize
\begin{enumerate}
  \item \textbf{Sentence graph construction:} split documents into sentences and create nodes; add edges for shared entities, semantic similarity, and sentence adjacency.
  \item \textbf{Question decomposition:} use the LLM to produce ordered sub-questions.
  \item \textbf{Seed retrieval:} retrieve candidate sentences for the current sub-question with embedding similarity and reranking.
  \item \textbf{Graph expansion:} expand from seed sentences to neighboring sentences so missing entity context can be recovered.
  \item \textbf{Answer and rewrite:} answer the current sub-question; if a later sub-question contains vague references, replace them with the recovered entity/answer.
  \item \textbf{Integration:} combine sub-answers or retrieved sub-contexts to answer the original question.
\end{enumerate}
\begin{block}{Core intuition}
ChainRAG treats multi-hop QA as a stateful chain: each answer repairs the next retrieval query before the system searches again.
\end{block}
\end{frame}

\begin{frame}{ChainRAG: Concrete Example}
\scriptsize
\begin{columns}[T,onlytextwidth]
  \begin{column}{0.48\textwidth}
    \begin{block}{Question}
      What was the hometown of the author of the novel about a male roe deer's life?
    \end{block}
    \begin{block}{Lost-in-retrieval}
      \textbf{Q1:} Who wrote the novel? $\rightarrow$ Felix Salten.\\
      \textbf{Bad Q2:} What was the home city of this author?\\
      The phrase ``this author'' loses the key entity.
    \end{block}
  \end{column}
  \begin{column}{0.48\textwidth}
    \begin{block}{ChainRAG repair}
      Rewrite the second query:
      \[
      \text{What was the home city of Felix Salten?}
      \]
      Retrieval now targets the correct entity instead of an unrelated author.
    \end{block}
  \end{column}
\end{columns}
\end{frame}

\begin{frame}{ChainRAG: Main Results}
\scriptsize
\begin{block}{Best F1 among ChainRAG variants in Table 1}
\centering
\resizebox{0.95\linewidth}{!}{
\begin{tabular}{lccc}
\toprule
LLM / Method & MuSiQue & 2Wiki & HotpotQA \\
\midrule
GPT4o-mini best baseline & 46.50 & 62.39 & 64.74 \\
GPT4o-mini ChainRAG & \textbf{50.54} & \textbf{62.55} & 64.59 \\
Qwen2.5-72B best baseline & 44.64 & 64.19 & 60.29 \\
Qwen2.5-72B ChainRAG & \textbf{49.37} & \textbf{65.85} & \textbf{64.54} \\
GLM-4-Plus ChainRAG & 51.66 & \textbf{70.58} & \textbf{64.22} \\
\bottomrule
\end{tabular}
}
\end{block}
\begin{block}{Retrieval repair}
Rewriting sub-question 2 raises Recall@2 to 58.81 on MuSiQue and 61.83 on HotpotQA, exceeding the original second-hop retrieval.
\end{block}
\end{frame}
